\subsection*{Bracket convention}
Throughout this article, brackets on
Lie algebras and Lie algebroids are defined using {\em
right}-invariant vector fields.

\section{Integrability}
\label{newsec}

In this section we prove Proposition \ref{frogP} by applying Crainic
and Fernandes' generalization of Lie III \cite{Crainic_Fernandes_03}.

Let $A$ be a Lie algebroid over $\Sigma$. In~\cite{Crainic_Fernandes_03, Crainic_Fernandes_11} the kernel of the
anchor of $A \rightarrow T \Sigma$ is denoted by~${\mathfrak g}
$. However, as this conflicts with our use as the Lie algebra of $G$,
we continue to denote the kernel by $A^\circ$. We otherwise follow the
notation and terminology of \cite{Crainic_Fernandes_11}.

In particular, the Weinstein groupoid of $A$ is denoted ${\mathcal
G}(A)$. An element of ${\mathcal G}(A)$ is a certain equivalence class
of $A$-paths. The obstruction to the existence of a bona fide Lie
groupoid integrating $A$ (that is, to the topological groupoid
${\mathcal G}(A)$ being a Lie groupoid) is measured by the {\it monodromy groups} $\tilde {\mathcal N}_{x_0}(A) $, $x_0 \in \Sigma
$. By definition, $\tilde {\mathcal N}_{x_0}(A) $ is the kernel of the
natural homomorphism
\begin{equation}
{\mathcal G}\big(A^\circ_{x_0}\big)\rightarrow {\mathcal G}(A)^\circ_{x_0} \label{dqp}.
\end{equation} On the right-hand side ${}^\circ$ denotes connected
component. At the level of Lie algebroid paths, this homomorphism is
just inclusion. The object on the left is a Lie group, while that on
the right may only be a topological group. Specializing
\cite{Crainic_Fernandes_03} to the transitive case, we have

\begin{Theorem}%\label{newsecT}
 Assuming its base manifold $\Sigma $ is connected, the Lie algebroid
 $A$ is integrable if and only if
 $\tilde {\mathcal N}_{x_0}(A) \subset {\mathcal G}\big(A^\circ_{x_0}\big)$
 is discrete, for some $x_0 \in \Sigma $.
\end{Theorem}

Now assume M1 and M2 hold. Then the restriction $\omega \colon
A^\circ_{x_0} \rightarrow {\mathfrak g} $ is an injection, integrating
to a homomorphism $\Omega \colon {\mathcal G}\big(A^\circ_{x_0}\big)
\rightarrow G$ {\em of Lie groups}, whose kernel $K_0$ is accordingly
discrete. On the other hand, we may include this homomorphism in the
following commutative diagram, whose vertical arrow is~\eqref{dqp}:
\begin{equation*}
 \begin{tikzcd} {\mathcal G}\big(A^\circ_{x_0}\big) \arrow{d}{}
\arrow{r}{\Omega} & G \\ {\mathcal G}(A)^\circ_{x_0}.
\arrow[swap]{ru}{} &
 \end{tikzcd}
\end{equation*} Here the diagonal map is the restriction of the natural
topological groupoid morphism ${\mathcal G}(A) \rightarrow {\mathcal
G}({\mathfrak g}) = G$, i.e., the map sending the equivalence class of
an $A$-path $a$ to the equivalence class of the ${\mathfrak g}$-path
$\omega \circ a$. Commutativity of the diagram implies the kernel of
the vertical map must lie in~$K_0$, but this kernel is, by definition,
$\tilde {\mathcal N}_{x_0}(A)$. This shows $\tilde {\mathcal
N}_{x_0}(A)$ is a discrete subset of~${\mathcal G}\big(A^\circ_{x_0}\big)$ and
the proof of Proposition~\ref{frogP}(1) now follows from the
theorem above.

Now let $x \in \Sigma $ be arbitrary and let ${\mathcal G} $ be a Lie
groupoid integrating $A$, which is necessarily transitive. Then,
assuming $\Sigma $ is connected, there exists an arrow $p \in
{\mathcal G} $ from $x_0$ to $x$. Conjugation $C_p\colon q \mapsto
pqp^{-1}$ maps the isotropy group ${\mathcal G}_{x_0}$ isomorphically
onto ${\mathcal G}_x$. Differentiating, we get a~Lie algebra
isomorphism ${\rm d}C_p \colon A^\circ_{x_0} \rightarrow A^\circ_x$ and a
commutative diagram
\begin{equation*}%\label{morph}
 \begin{CD} A^\circ_{x_0} @>{{\rm d}C_p}>> A^\circ_x \\ @V{\omega}VV
@VV{\omega}V \\ {\mathfrak g} @>{\Adjoint_{\Omega(p)}}>> {\mathfrak g}.
 \end{CD}
\end{equation*}
From this observation parts~(2) and~(3) of Proposition \ref{frogP} immediately follow.

\section{The uniqueness of primitives up to symmetry}\label{infch}
This section establishes the uniqueness of primitives, appropriately
defined, up to symmetry. The central result is Theorem
\ref{uniquenessT}. Combining our uniqueness result with Theorem
\ref{outlineT}, we obtain the existence and uniqueness Theorem
\ref{summaryT}.

\subsection*{Symmetries of a homogeneous space}
% \label{symmetries}

For the purposes of constructing a theory with the proper invariance,
we have been regarding our fixed data as a homogeneous
$G$-space. While a choice of point $m_0 \in M$ gives us an
identification $M \cong G/G_{m_0}$, formulations depending on a choice
of base point are to be eschewed.\footnote{Geometries in the real
 world do not come with a preferred choice of base-point. Base-points
 are an artifact of Klein's abstraction of geometry, not an intrinsic
 feature.} This decision has a somewhat unexpected consequence,
anticipated by reconsidering the simplest case.

According to Cartan's theorem, a smooth map $f \colon \Sigma
\rightarrow G$ is uniquely determined by its logarithmic derivative,
up to symmetries of $G$. Here a ``symmetry'' is a right group
translation. However, to obtain an invariant version of Cartan's
result we must broaden both the notion of symmetry and what it means
to be a primitive. To see why, consider a smooth map $f \colon \Sigma
\rightarrow M$, where $ M$ is a smooth manifold on which $G$ is acting
transitively and {\em freely}, so that $M \cong G$, up to choice of
base-point. In~order to drop the right-invariant Maurer--Cartan form
$\omega_G$ on $G$ to a~one-form on $ M$, we must suppose here that $G$
is acting on $ M$ from the {\em left}. For then, fixing $ m_0 \in M$
and defining a $\Phi(g) = g \cdot m_0 $, the diffeomorphism $\Phi
\colon G \rightarrow M$ pushes $\omega_G$ forward to a~one-form
$\omega_M $ on $M$ that is independent of the choice of $ m_0$. But
then $\omega_M$ is {\em not} invariant with respect to the action of
$G$~--- rather it is {\em equivariant}, if we regard $G$ as acting on
${\mathfrak g} $ by adjoint action. In~particular, two smooth maps
$f_1, f_2 \colon \Sigma \rightarrow M$ with $f_2(x)=g \cdot f_1(x)$,
$x \in \Sigma $, have, in~general, {\em different} logarithmic
derivatives: $f_2^*\omega_M = \langle\Adjoint_g,f_1^* \omega_M\rangle
$.

To proceed one defines $f \colon \Sigma \rightarrow M$ to be a {\it primitive} of a one-form $\omega $ on~$\Sigma $ if $f^*\omega_M $ and~$\omega $ agree ``up to adjoint action''. The price one pays for this
relaxed definition is that the logarithmic derivative $f^* \omega_M $
only determines $f$ up to a larger class of symmetries of $M$. Under
an identification $M \cong G$, these symmetries consist of the
diffeomorphisms generated by all right {\em and} left translations.

Symmetries in the general case are formalised as follows:
\begin{Definition}\label{symmetriesD}
 Let $M$ be a homogeneous $G$-space. Then a {\it symmetry} of $M$ is
 any diffeomorphism $\phi \colon M \rightarrow M$ for which there
 exists some $l \in G$ such that
 $\phi(g \cdot m) = lgl^{-1} \cdot \phi(m)$ for all
 $g \in G$, $m \in M$.
\end{Definition}

 The symmetries of $M$ form a Lie group henceforth denoted
$\automorphism(M)$. Evidently, $\automorphism(M)$ contains every left
translation $\phi(m):=k \cdot m$, $k \in G$ (take $l=k$). The
following characterization of symmetries is readily verified:

\begin{Proposition}\label{symmetriesP}
 Fix a point $m_0 \in M$ and identify $M$ with the left coset space
 $G/H$, where $H$ denotes the isotropy at $m_0$. Let $l \in G$ be
 arbitrary and suppose $r \in G$ is in the normaliser of $H$, so that
 there exists a map $\phi \colon G/H \rightarrow G/H$ making the
 following diagram commute:
 \begin{equation*}
 \begin{CD} G @>{g~\mapsto~lgr^{-1}}>> G \\ @VVV @VVV \\ G/H
@>{\phi}>> G/H.
 \end{CD}
 \end{equation*}
 Then $\phi $ is a symmetry of $M$ and all
symmetries of $M$ arise in this way. In~other words, the Lie group
$W:=N_G(H)/H$ acts on the left of $G/H$ according to
\[
rH \cdot gH = gr^{-1}H,
\]
$($an action commuting with the left action of $G)$ and
$\automorphism(M)$ is the Lie group generated by both left
translations and those transformations of $M \cong G/H$ defined by the
action of $W$. Here $N_G(H)$ denotes the normaliser of $H $ in $G$.
\end{Proposition}
